\section{Introduction}
\label{sec:intro}

Voice interfaces can process audio remotely or locally. Local processing can reduce network dependence and keep audio on the device, but requires computation and storage within its power and area budget. For a model fixed over a product's lifetime, mask-programmed weights offer a way to avoid off-chip fixed-weight transfers. Local readout, activation storage and arithmetic still consume energy.

Hardwired, or ``model-specific'', inference silicon has recently attracted commercial attention at the frontier, where fixed-weight chips promise very high token rates for large language models~\cite{taalas2026hc1}. This paper asks the complementary question at the other end of the scale:
\begin{quote}
\emph{Can a useful speech model be hardwired into an open, mature, low-cost process, and can a very small team build and verify it using AI coding agents and open-source tools?}
\end{quote}

We present RouteVoice, a Moonshine Tiny speech recognizer~\cite{jeffries2024moonshine} designed for GlobalFoundries' open 180\,nm process (GF180MCU). Eligible matrices use fixed 4-bit indices into local codebooks. Their integer projection products are assembled from shared partials; frontend INT8 storage, FP32 operations and dynamic attention remain.

The building blocks are not new:
\begin{itemize}
\item codebook-indexed weights~\cite{han2016eie,zadeh2022mokey};
\item shared multiples of an input selected per fixed coefficient~\cite{park2002cshm}, which hardwired-LLM patents also describe~\cite{taalas2025patent802,taalas2026patent566};
\item hardwired weights on an open PDK~\cite{pai2026ankhdjet}.
\end{itemize}
What we add is their combination in a speech recognizer verified at RTL, plus a separately routed projection engine and model-derived index banks. Full-chip physical integration remains incomplete.

\paragraph{Contributions.}
\begin{enumerate}
\item \textbf{A hardwired-codebook weight format, validated for ASR.} K64 stores each weight as a 4-bit index into a 16-entry UINT8 book shared by 64 weights. It keeps Moonshine Tiny within $+2$ WER points of the original on a reserved test, and cuts storage of seven converted RTL services by 39\% (not total model storage).
\item \textbf{A bit-exact recognizer.} The full recognizer RTL decodes speech to the exact tokens and intermediate hashes of the frozen software model, with codebook services, an affine INT8 frontend and unchanged supporting tables.
\item \textbf{A routed, signoff-checked projection engine.} The $1280\times320$ engine (409,600 MACs per projection) closes timing at \SI{150}{ns} in all nine PVT/RC corners. Routing, Magic DRC, KLayout DRC, antenna and LVS checks are clean. Estimated activity-annotated routed energy is \SI{287}{pJ}/MAC; slew, capacitance and fanout violations remain.
\item \textbf{Full-design evidence.} Composed per-module synthesis for silicon and FPGA completes with weight ROMs black-boxed. Mask-programmed banks passing selected DRC checks cover all 98 matrices (\SI{174.7}{mm^2}), with full LVS on six representative banks.
\item \textbf{A documented agent-driven methodology.} We report the step-by-step procedure, the failures and their fixes, and the cost (\$34.82 of cloud compute plus a laptop), and derive concrete practices for faster flows.
\item \textbf{An analysis of future scope.} We size what the same approach implies for local LLMs, and propose ways to make hardwired silicon serve multiple tasks.
\end{enumerate}

\Cref{sec:background} reviews related work. \Cref{sec:representation,sec:architecture} describe the weight format and architecture, and \cref{sec:method} the methodology. \Cref{sec:results,sec:closure} report results, including a timing-closure case study. \Cref{sec:efficiency,sec:future,sec:limitations} discuss efficiency, future scope and limitations.
